\chapter{Background}

\section{Metric First Order Temporal Logic}
% TODO: introduce (E)MFOTL

\section{The Lex Language}

\subsection{Constitutive, Regulative, and Exception Rules (CRE Meta-Logic)}

\section{Proactive Enforcement}

First we recall some definitions from Metric First Order Temporal Logic (MFOTL) and Lex \cite{whyenf, lex}.
Let $\Types$ be a set of base types and $\Variables$ a set of variables.

\begin{definition}\label{def:signature}
    A \emph{Signature}
    $\Sigma = \left( \left(\Domain_{ty}\right)_{ty \in \Types}, \Events, \Functions, \iota \right)$
    is a quadruple consisting of a domain of values such that
    for every $ty \in \Types$, $\Domain_{ty}$ is the set of values of type $ty$,
    a set of event names $\Events$, a set of function names $\Functions$,
    and an arity function $\iota : \Events \cup \Functions \to \Types^*$
    such that for every $e \in \Events$, $\iota(e) = (p_1,\dots,p_k)$ provides the
    parameter types of $e$
    and for any $f \in \Functions$, $\iota(f) = (a_1,\dots,a_k,r)$ provides the
    argument types $(a_1,\dots,a_k)$ and the return type $r$ of $f$.
\end{definition}

\begin{definition}\label{def:event}
    An \emph{Event} $e(p_1,\dots,p_k)$ consists of an event name
    $e \in \Events$ and parameters $(p_1,\dots,p_k)$ such that
    for every $1 \leq i \leq k$, $p_i \in \Domain_{\iota(e)_i}$.
\end{definition}

\begin{definition}\label{def:database}
    A \emph{Database} $D \in \Databases$ is a set of events.
\end{definition}

\begin{definition}\label{def:trace}
    A \emph{Trace} $\sigma \in \Traces$ is a sequence
    $\Sequence[1 \leq i \leq k]{(\tau_i, D_i)}$, $k \in \Nat \cup \lbrace \infty \rbrace$
    of timestamps $\tau_i \in \Nat$ and finite databases $D_i \in Databases$,
    where timestamps grow \emph{monotonically}
    ($\forall i < |\sigma|.\,\tau_i \leq \tau_{i+1}$)
    and \emph{progress}
    (if $|\sigma| = \infty$, then $\lim_{i\to\infty} \tau_i = \infty$).
\end{definition}

\begin{definition}\label{def:infinite-trace}
    The set of \emph{Infinite Traces} is denoted by $\Traces_\omega$.
\end{definition}

\begin{definition}\label{def:finite-trace}
    The set of \emph{Finite Traces} is denoted by $\Traces_f$.
\end{definition}

\begin{definition}\label{def:time-point}
    A \emph{time-point} $0 \leq i \leq |\sigma|$ is an index in a trace $\sigma$.
\end{definition}

\begin{definition}\label{def:empty-trace}
    The \emph{Empty Trace} is denoted by $\emptytrace$.
\end{definition}

\begin{definition}\label{def:property}
    A \emph{Property} is a subset $P\subseteq \Traces_\omega$.
\end{definition}

A Property can be induced by some policy $\phi$, we denote this policy by
$P(\phi)$.
We say a trace $\sigma$ is compliant with policy $\phi$ if and only if
$\sigma \in P(\phi)$.

% TODO: introduce the notion of a System under Scrutiny (SuS)

\begin{definition}
    \label{def:enf-function}
   An \emph{Enforcer Function}
   $\mu_\phi : \Traces \times \Nat \times (\Databases \cup \lbrace\tick\rbrace) \to (\Databases \cup \lbrace \bot \rbrace)$
   for some policy $\phi$ takes as input a trace $\sigma \in \Traces$,
   a timestamp $\tau \in \Nat$, and either a database $D \in \Databases$
   or the special event $\tick$ which symbolizes a clock tick of the underlying
   System under Scrutiny (SuS) and outputs either
   a database $D' \in \Databases$ or nothing (marked by the special symbol $\bot$)
   and if the trace $\sigma \in P(\phi)$, it ensures that
   $\sigma \cdot \left(\tau, D'\right) \in P(\phi)$.
   $\bot$ (nothing) is only a valid output if the third argument (the event) is $\tick$.
   % TODO: is this remark about \bot / nothing correct?
\end{definition} % TODO: is there an additional requirement that \tau must be \geq than the largest timestamp in \sigma?

% TODO: elaborate on functionality of a proactive enforcer

Capabilities are an integral part of the proactive enforcement
approach introduced by Hublet et. al \cite{whyenf}.
Under this approach an enforcer can have different capabilities, or ways of interacting
with events emitted by a given system under scrutiny (SuS).
In Hublet et. al the set of capabilities consists of observable, suppressable, causable,
and suppressable-causable.
Hublet et. al \cite{lex} expands this expands this system of capabilities, specifically it
makes the existence of events that may not (always) be externally visible explicit.
Under both notions a capability function $\kappa : \Events \to \Capabilities$ is used
to assign a capability to all events with some event name $e \in \Events$.
Also in both contexts, capabilities form a (semi)-lattice $(\Capabilities, \sqsubseteq)$.
We write Sup as a shortform for Suppressable and Cau as a shortform for Causable.
% TODO: insert WhyEnf and Lex/CRE capability lattices

\begin{definition}
    \label{def:cau-events}
    The subset of \emph{Causable Events}
    $\Events_{Cau} = \left\lbrace e~|~\text{Cau} \sqsubseteq \kappa(e) \right\rbrace$
    consists of all event names which are can be caused by the enforcer.
\end{definition}

\begin{definition}
    \label{def:cau-events}
    The subset of \emph{Suppressable Events}
    $\Events_{Sup} = \left\lbrace e~|~\text{Sup} \sqsubseteq \kappa(e) \right\rbrace$
    consists of all event names which can be suppressed by the enforcer.
\end{definition}

\begin{theorem}
    A Function $\mu_\phi : \Traces \times \Nat \times \Databases \to \Databases$
    is a valid enforcer function for some (Lex or MFOTL) policy $\phi$ iff
    for all $\sigma \in \Traces$, $\tau \in \Nat$, and $D \in \Databases$
    \begin{enumerate}
        \item $\mu_\phi(\sigma, \tau, D) \setminus D \subseteq \Events_{Cau}$ and
        \item $D \setminus \mu_\phi(\sigma, \tau, D) \subseteq \Events_{Sup}$
    \end{enumerate}
\end{theorem}

\begin{proof}
    % TODO
    TODO
\end{proof}

\section{Refinement Mappings}
Abadi and Lamport \cite{refinement-mappings}

TODO

\section{Enforcement as an Event System}
% TODO: probably move out of background chapter

We first define enforcement as an event system and then show how the notion of refinement as in Sprenger et. al \cite{igloo} works for such an event system.

\begin{definition}
    \label{def:enf-state-space}
    The \emph{Enforcement State Space} $\Statespace = \Nat \times \Traces$
    contains tuples $(\tau,\sigma)$.
\end{definition}
The timestamp $\tau$ can be seen as the \emph{current} timestamp
and $\sigma$ the trace up until timestamp $\tau$.

\begin{definition}
    \label{def:enf-initial-states}
    The set of \emph{Initial States}
    $\Initialstates = \lbrace (\tau, \emptytrace) | \tau \in \Nat \rbrace$
    consists of all states with an arbitrary timestamp $\tau$ and the empty trace
    $\emptytrace$.
\end{definition}
Intuitively this can be interpreted as that an enforced system may be
started at an arbitrary timestamp.

\begin{definition}
    \label{def:enf-event-space}
    The \emph{Event Space} $\Eventspace=\Databases \cup \lbrace\tick\rbrace$ is the union of databases and the special event $\tick$ which stands for a clock tick
    (incrementing a timestamp by one).
\end{definition}

\begin{definition}
    \label{def:enf-transition-relation}
    The \emph{Transition Relation} is a subset
    $\Transitionrelation{\phi}{} \subseteq \Statespace \times \Eventspace \times \Statespace$.
    We write $(\tau, \sigma) \Transitionrelation{\phi}{e} (\tau', \sigma')$
    as a shorthand for
    $\left(\left(\tau,\sigma\right), e, \left(\tau',\sigma'\right)\right) \in \Transitionrelation{\phi}{}$.
    A Transition Relation is valid for some policy $\phi$, if there exists an
    enforcer function $\mu_\phi$ for $\phi$ and all allowed transitions match
    one of the three following transitions.
    \begin{itemize}
        \item $(\tau, \sigma) \Transitionrelation{\phi}{\mu_\phi(\sigma, \tau, D))} (\tau, \sigma \cdot (\tau, \mu_\phi(\sigma, \tau, D))$
        \item $(\tau, \sigma) \Transitionrelation{\phi}{\mu_\phi(\sigma, \tau, \tick))} (\tau, \sigma \cdot (\tau, \mu_\phi(\sigma, \tau, \tick))$ if $\mu_\phi(\sigma, \tau, \tick) \neq \bot$
        \item $(\tau, \sigma) \Transitionrelation{\phi}{\mu_\phi(\sigma, \tau, \tick))} (\tau, \sigma)$ if $\mu_\phi(\sigma, \tau, \tick) = \bot$
    \end{itemize}
\end{definition}

\subsection{Refinement of Enforcement Event Systems}

For Enforcement Event Systems we will now define refinement
via a modulo function and a simulation relation as in Sprenger et. al \cite{igloo}.
We aim to show that for a specific class of mediator functions, whenever
some concrete enforcement event system refines an abstract enforcement 
event system, the existence of an enforcer function $\mu_A$ for the abstract
enforcement event system implies the existence of an enforcer function $\mu_C$ for
the concrete enforcement event system.
We will further show that the way we refine the enforcement event systems here
corresponds to the refinement definition for Lex in Hublet et. al \cite{lex}.

We start out with two enforcement event systems, an abstract and a concrete one.
The parameters of the abstract system are the following
\begin{align*}
    \Sigma_A &= \left(\Domain, \Events_A, \Functions, \iota_A \right) \\
    \Statespace_A &= \Nat \times \Traces_A & \text{State space of traces of signature $\Sigma_A$}\\
    \Initialstates_A &= \left\lbrace \left(\tau, \emptytrace \right) | \tau \in \Nat \right\rbrace \\
    \phi_A &= \left( C_A, R_A, E_A \right) &\text{An enforceable CRE policy} \\
    S_A &= \left(\Sigma_A, \phi_A, \kappa_A \right)
    % \\
    % \mu_{\phi_A} && \text{An enforcer function for policy $\phi_A$} \\
\end{align*}

And the parameters for the concrete system are
\begin{align*}
    \Sigma_C &= \left(\Domain, \Events_C, \Functions, \iota_C \right) \\
    \Statespace_C &= \Nat \times \Traces_C & \text{State space of traces of signature $\Sigma_C$}\\
    \Initialstates_C &= \left\lbrace \left(\tau, \emptytrace \right) | \tau \in \Nat \right\rbrace \\
    \phi_C &= \left( C_C, R_C, E_C \right) \\
    S_C &= \left(\Sigma_C, \phi_C, \kappa_C \right)
    % \\
    % \mu_{\phi_C} && \text{An enforcer function for policy $\phi_C$} \\
\end{align*}

We assume that the concrete specification $S_C$ is a refinement of the
abstract specification $S_A$ as defined in Hublet et. al \cite{lex}.

\begin{definition}
    \label{def:modulo-mediator-function}
    A \emph{Modulo Mediator Function}
    $\pi : \Traces_C \times \Nat \times \Databases_C \to \Databases_A$
    converts a set of concrete events $D_C$ to a set of abstract events
    $D_A = \pi(\sigma_C, \tau, D_C)$.
\end{definition}

This definition of a modulo mediator function differs from the definition in
Sprenger et. al \cite{igloo} by the inclusion of a trace and timestamp in its arguments.

\begin{definition}
    \label{def:simulation-relation}
    A \emph{Simulation Relation}
    $R \subseteq \Statespace_C \times \Statespace_A$ defines a correspondence of
    concrete and abstract states.
    We define $R$ with respect to some modulo mediator function $\pi$ as
    $$
    R = \left\lbrace \left( \left(\tau, \sigma_C \right), \left(\tau, \sigma_A \right) \right) | \left|\sigma_C\right| = \left|\sigma_A\right| \wedge \forall i \leq \left| \sigma_C \right| .\, \pi\left( \sigma_{C,\dots<i}, \tau, \sigma_{C,i} \right) = \sigma_{A,i} \right\rbrace
    $$
\end{definition}
Where $\sigma_{C,\dots<i}$ is the trace $\sigma$ with all time-points less than $i$.

\begin{theorem}
    \label{theo:lex-refinement-and-event-system-refinement}
    If a concrete event system is a refinement of the abstract event system
    through the modulo mediator function $\pi$ and the simulation relation $R$,
    then the concrete specification $S_C$ is a refinement of the abstract
    specification $S_A$.
\end{theorem}

\begin{proof}
    \label{proof:refinement-correspondence}
    TODO
\end{proof}

\begin{theorem}
    \label{theo:existence-of-concrete-enforcer}
    If there exists and abstract enforcer $\mu_A$ and a concrete event system refines
    the event system defined by $\mu_A$, then there exists a concrete enforcer
    $\mu_C$.
\end{theorem}

\begin{proof}
    \label{proof:existence-proof}
    Let $\phi_A$ be some enforceable policy for which there exists an enforcer
    function $\mu_A$ and
    \begin{align}
        \label{ex-proof:c-states}
        \left( \tau, \sigma_C \right), \left( \tau', \sigma_C'  \right) &\in \Statespace_C \\
        \label{ex-proof:a-states}
        \left( \tau, \sigma_A \right) &\in \Statespace_A \\
        \label{ex-proof:c-a-state-in-r}
        \left(\left( \tau, \sigma_C \right), \left( \tau, \sigma_A \right) \right) &\in R \\
        \label{ex-proof:e}
        e &\in \Eventspace \\
        \label{ex-proof:concrete-transition} 
        \left( \tau, \sigma_C \right) \Transitionrelation{\phi_C}{e} \left( \tau', \sigma_C'  \right) \\
    \end{align}
    Now we will do a case distinction to show that for all three cases
    of transitions in Definition \ref{def:enf-transition-relation}
    there must exist an abstract state $\left( \tau', \sigma_A' \right) \in \Statespace_A$
    such that 
    $$
        \left( \tau, \sigma_A \right) \Transitionrelation{\phi_C}{\pi\left(\sigma_C, \tau, e\right)} \left( \tau', \sigma_C'  \right) \\
    $$
\end{proof}